
et par
suite, pour tout @@M0@@ :
@@M1@@

D'après (7.1.1), les @@M2@@ sont des polynômes isobares de poids @@M3@@ par rapport aux @@M4@@, avec @@M5@@. Donnons par exemple les valeurs de @@M6@@ et @@M7@@ :
@@M8@@

7.2. Soit @@M9@@ un anneau de caractéristique @@M10@@~; supposons données des applications @@M11@@ de @@M12@@ dans @@M13@@, additives, et telles que @@M14@@. On définit alors une pré-@@M15@@-structure sur @@M16@@ en posant :
@@M17@@
La multiplicativité résulte de l'additivité des @@M18@@. Il y a donc équivalence entre la donnée d'une pré-@@M19@@-structure sur @@M20@@, et la donnée des opérations @@M21@@ correspondantes.

7.3. Il est clair d'après la formule (7.1.1) que les @@M22@@ commutent aux @@M23@@-homomorphismes. Si @@M24@@ sont des pré-@@M25@@-anneaux de caractéristique nulle, il est nécessaire et suffisant qu'un homomorphisme @@M26@@ commute aux @@M27@@ pour qu'il soit un @@M28@@-homomorphisme~: cela résulte de 7.2, les @@M29@@ s'expriment comme polynômes par rapport aux @@M30@@.

\textbf{Exemples 7.4.} i) Pour @@M31@@, muni de sa structure de @@M32@@-anneau, on a :
@@M33@@
Par suite, les @@M34@@ sont tous égaux à l'application identique de @@M35@@.

ii) Plaçons nous sur la catégorie @@M36@@ des @@M37@@-algèbres (1.8). Les 
